\chapter{Normalformen und Darstellungen von Halbverbänden}
\label{b45guide:ch:semilattices}

\section{Terme und Erzeugung}

\begin{definition}[Halbverbandsterm und Variablenträger]
Ein \emph{Halbverbandsterm} ist ein endlicher Ausdruck aus Variablen
und einer binären Verknüpfung $\vee$, ohne Konstanten.
Sein \emph{Variablenträger} $\operatorname{Var}(t)$ ist die Menge
der in ihm auftretenden Variablen; sie ist endlich und nichtleer.
Eine \emph{Identität} $s=t$ gilt in einem Halbverband, wenn beide
Terme bei jeder Belegung ihrer Variablen denselben Wert haben.
\end{definition}

\begin{theorem}[Normalform und vollständiges Identitätskriterium]
\label{b45guide:thm:support}
In jedem Halbverband hängt der Wert eines Terms nur von seinem
Variablenträger ab. Reihenfolge, Klammerung und positive
Wiederholungszahlen der Variablen ändern den Wert nicht.
Für zwei Halbverbandsterme $s,t$ gilt
\[
 \bigl(s=t\text{ gilt in allen Halbverbänden}\bigr)
 \quad\Longleftrightarrow\quad
 \operatorname{Var}(s)=\operatorname{Var}(t).
\]
Jede Identität, die nicht allgemein gilt, scheitert bereits in
dem zweielementigen Halbverband $\{0,1\}$ mit
$0\vee0=0$ und $0\vee1=1\vee0=1\vee1=1$. \cite{smith}
\end{theorem}

\begin{definition}[Endliche Verknüpfung und erzeugter Halbverband]
Für eine endliche nichtleere Teilmenge $F$ eines Halbverbandes
$(S,\vee)$ bezeichnet $\bigvee F$ die Verknüpfung ihrer Elemente;
die Normalform macht sie unabhängig von Aufzählung und Klammerung.
Eine Teilmenge $U\subseteq S$ heißt \emph{Unterhalbverband}, wenn
$x\vee y\in U$ für alle $x,y\in U$ gilt.
Der von $X\subseteq S$ \emph{erzeugte Unterhalbverband}
$\langle X\rangle_{\vee}$ ist der kleinste Unterhalbverband, der $X$ enthält.
\end{definition}

\begin{theorem}[Explizite Erzeugung und lokale Endlichkeit]
\label{b45guide:thm:semigen}
Für jeden Halbverband $(S,\vee)$ und jedes $X\subseteq S$ gilt
\[
 \langle X\rangle_{\vee}
   =\{\bigvee F\mid \varnothing\ne F\subseteq X,\ F\text{ endlich}\}.
\]
Insbesondere ist $\langle\varnothing\rangle_{\vee}=\varnothing$.
Ist $|X|=n<\infty$, so
\[
 |\langle X\rangle_{\vee}|\leq 2^n-1.
\]
Jeder endlich erzeugte Halbverband ist daher endlich. Die Schranke
wird für jedes $n\geq0$ erreicht. \cite{smith,supclosed}
\end{theorem}

\section{Freie Halbverbände}

\begin{definition}[Freie Erzeugung]
\label{b45guide:def:free-semilattice}
Ein Halbverband $F$ mit einer Abbildung $\iota:X\to F$ heißt
\emph{frei auf $X$}, wenn für jeden Halbverband $T$ und jede
Abbildung $g:X\to T$ genau eine Abbildung $\widehat g:F\to T$
existiert, die
\[
 \widehat g(u\vee v)=\widehat g(u)\vee\widehat g(v),
 \qquad \widehat g\circ\iota=g
\]
erfüllt. Solche operationserhaltenden Abbildungen heißen
Homomorphismen; Kapitel~\ref{b45guide:ch:homomorphisms}
präzisiert ihre Eigenschaften.
\end{definition}

\begin{theorem}[Freier Halbverband als endliche Mengenfamilie]
\label{b45guide:thm:free-semilattice}
Für jede Menge $X$ ist
\[
 \mathcal P_{\mathrm{fin}}^{+}(X)
  =\{F\subseteq X\mid F\text{ endlich und }F\ne\varnothing\}
\]
mit Vereinigung und $\iota(x)=\{x\}$ frei auf $X$.
Die eindeutige Fortsetzung von $g:X\to T$ ist
\[
 \widehat g(F)=\bigvee_{x\in F}g(x).
\]
Freie Halbverbände auf denselben benannten Erzeugern sind eindeutig
bis auf einen diese Erzeuger festlassenden Isomorphismus.
Bei $|X|=n<\infty$ besitzt der freie Halbverband genau $2^n-1$ Elemente.
Für $X=\varnothing$ ist er gemäß der Trägerkonvention leer. \cite{smith}
\end{theorem}

\begin{theorem}[Freier Halbverband mit Null]
\label{b45guide:thm:free-zero}
Werden eine Null und ihre Erhaltung durch Homomorphismen verlangt,
so ist der freie Halbverband auf $X$ die Struktur
\[
 (\mathcal P_{\mathrm{fin}}(X),\cup,\varnothing).
\]
Es gilt $\widehat g(\varnothing)=0$; auf nichtleeren Mengen gilt
die Formel aus Theorem~\ref{b45guide:thm:free-semilattice}.
Bei $|X|=n<\infty$ beträgt seine Größe $2^n$.
In einem Halbverband mit Null ist entsprechend die unter
$\vee$ und $0$ erzeugte Teilstruktur
\[
 \langle X\rangle_{\vee,0}
   =\{0\}\cup\langle X\rangle_{\vee}.
\]
Für endliche Verknüpfungen wird dann $\bigvee\varnothing=0$ gesetzt.
\cite{smith}
\end{theorem}

\section{Darstellung durch Mengen und zweielementige Faktoren}

\begin{theorem}[Vereinigungsdarstellung jedes Halbverbandes]
\label{b45guide:thm:union-representation}
Sei $(S,\vee)$ ein beliebiger Halbverband mit seinem induzierten
Vergleich. Für $x\in S$ setze
\[
 \Phi(x)=\{a\in S\mid x\not\leq a\}.
\]
Dann gilt
\[
 \Phi(x\vee y)=\Phi(x)\cup\Phi(y),\qquad
 \Phi(x)\subseteq\Phi(y)\iff x\leq y .
\]
Damit ist $\Phi$ ein Isomorphismus von $(S,\vee)$ auf die
unter Vereinigung abgeschlossene Familie $\Phi[S]\subseteq\mathcal P(S)$.
Besitzt $S$ eine Null, so ist $\Phi(0)=\varnothing$.
\end{theorem}

\begin{theorem}[Durchschnittsdarstellung der infimalen Lesart]
\label{b45guide:thm:intersection-representation}
Sei $(S,\wedge)$ infimal gelesen. Dann liefert
\[
 \Psi(x)=\{a\in S\mid a\leq x\}
\]
einen Isomorphismus auf eine unter Durchschnitt abgeschlossene
Mengenfamilie:
\[
 \Psi(x\wedge y)=\Psi(x)\cap\Psi(y).
\]
Im Falle eines Verbandes ist diese Darstellung zunächst eine
Darstellung seines $\wedge$-Halbverbandes. Eine gleichzeitige
Darstellung beider Verbandsoperationen als gewöhnlicher
Durchschnitt und gewöhnliche Vereinigung folgt daraus nicht.
\end{theorem}

\begin{theorem}[Trennung durch zweielementige Halbverbände]
\label{b45guide:thm:two-separation}
Sei $\mathbf2$ der zweielementige Halbverband aus
Theorem~\ref{b45guide:thm:support}. Zu jedem $a\in S$ ist
\[
 h_a(x)=
 \begin{cases}
  0,&x\leq a,\\
  1,&x\not\leq a
 \end{cases}
\]
ein Halbverbandshomomorphismus $S\to\mathbf2$.
Für $x\ne y$ existiert ein $a\in S$ mit $h_a(x)\ne h_a(y)$.
Folglich ist
\[
 x\longmapsto(h_a(x))_{a\in S}
\]
eine Einbettung in den Halbverband $\mathbf2^S$ mit
koordinatenweiser Verknüpfung. Auch bei unendlichem $S$
werden nur zweielementige Faktoren benötigt.
\end{theorem}
